\section{Standards of Evidence}

The classification, mathematical forms, and normative claims should be assessed separately.

A classification earns its place by helping identify mechanisms or distinguish cases. A quantitative model needs defined variables, stable boundaries, calibrated parameters, and out-of-sample evaluation. A normative proposal needs public reasons about interests, rights, authority, and distribution.

Physical terms require a demonstrated correspondence when applied elsewhere. A biological thickening model needs an input, a response, and a mechanism. A social Deborah number needs an identifiable adjustment process and timescale. Similar language or curve shape supplies a hypothesis to investigate.

The exponential form should compete with interaction, threshold, saturation, and simpler alternatives. If adequate fit requires repeatedly changing what the coordinates mean, revisit the definitions. If several media show the same curve, inspect the mechanism before deciding what the similarity implies.

The framework's central practical claim is that proposals become dependable changes through a chain with receiving capacities, feedback, and obligations. Its usefulness can be tested by whether following that chain improves decisions.

The book's normative commitments are explicit: affected people need effective participation and remedies; consequential judgments need evidence; authority needs an accountable scope; and progress should be assessed over the relevant system and time.

Application begins with feasible paths and a full account of cost. It then considers the interests and rights involved, establishes authority and responsibility, and preserves records through which the result can continue to be challenged and improved.
